% Full proofs for the exact propositional LJ / Kripke consequence presentation.
\begin{annexitem}{annex:p8.ipl.complete.persistence}{thm:p08-ipl-completeness-persistence}{Persistence for all propositional formulas}
\begin{proof}
Induct on the formation of \(A\), with arbitrary worlds \(w\leq v\).
For an atom, the assertion is the persistence assumption on \(V\).
For \(\bottom\), the antecedent of the assertion is impossible.
If \(w\forces B\land C\), then \(w\) forces both conjuncts.
Each persists to \(v\) by induction, so \(v\forces B\land C\).
If \(w\forces B\lor C\), one of the disjuncts is forced at \(w\).
That disjunct persists to \(v\), so \(v\forces B\lor C\).

Finally, suppose \(w\forces B\to C\). To show
\(v\forces B\to C\), take any \(u\geq v\) forcing \(B\).
Transitivity gives \(u\geq w\), and the implication clause at \(w\)
gives \(u\forces C\). This proves the implication case. Applying the
result separately to each member of a premise set proves persistence for
that set, without requiring it to be finite.
\end{proof}
\end{annexitem}

\begin{annexitem}{annex:p8.ipl.complete.soundness}{thm:p08-ipl-completeness-soundness}{Soundness of the thirteen propositional LJ rule schemes}
\begin{proof}
A finite-multiset sequent is valid when, at every world of every model,
forcing each formula occurrence in its antecedent implies forcing its
single succedent. Induct on the finite derivation tree. For each rule,
assume its premise sequents are valid, fix a model, and let \(w\) force
the entire conclusion antecedent. The following cases cover the exact
propositional sheet.

\paragraph{Id and falsum left.}
The identity sequent \(B\Rightarrow B\) is valid by its assumption.
The antecedent of \(\bottom,\Gamma_0\Rightarrow C\) cannot be forced
at any world, because no world forces \(\bottom\). This proves
\(\bottom L\).

\paragraph{Weakening and contraction.}
For WL, forcing \(B,\Gamma_0\) includes forcing \(\Gamma_0\), so
the premise gives its succedent. For CL, a world forcing
\(B,\Gamma_0\) forces both antecedent occurrences of \(B\) in the
premise \(B,B,\Gamma_0\Rightarrow C\); validity of that premise
gives \(C\). Multiplicity records proof resources here, but the forcing
condition on each occurrence is the same condition.

\paragraph{Conjunction left and right.}
For \(\land L\), forcing \(B\land C,\Gamma_0\) gives forcing of
\(B,C,\Gamma_0\); the premise gives the succedent. For
\(\land R\), the two premise sequents have the same antecedent
\(\Gamma_0\). Their validity gives both conjuncts at \(w\), and
therefore gives their conjunction at \(w\).

\paragraph{Disjunction left and both right rules.}
For \(\lor L\), forcing \(B\lor C,\Gamma_0\) means either
\(w\forces B\) or \(w\forces C\). In the first case use the premise
\(B,\Gamma_0\Rightarrow D\); in the second use
\(C,\Gamma_0\Rightarrow D\). Each gives \(w\forces D\).
For \(\lor R_1\), the premise gives the left disjunct at \(w\),
which suffices for the conclusion disjunction. For \(\lor R_2\),
the premise gives its right disjunct instead.

\paragraph{Implication left.}
The premises are \(\Gamma_0\Rightarrow B\) and
\(C,\Pi_0\Rightarrow D\); the conclusion antecedent is
\(B\to C,\Gamma_0,\Pi_0\). The first premise gives
\(w\forces B\). Reflexivity makes \(w\) one of its own future
worlds, so \(w\forces B\to C\) gives \(w\forces C\).
Together with the already forced \(\Pi_0\), this supplies the entire
antecedent of the second premise and hence \(w\forces D\).
The two split contexts need not be equal or disjoint.

\paragraph{Implication right.}
The premise is \(B,\Gamma_0\Rightarrow C\). For any
\(v\geq w\) with \(v\forces B\), persistence gives
\(v\forces\Gamma_0\) from \(w\forces\Gamma_0\). Validity of the
premise at \(v\) now gives \(v\forces C\). This holds at every such
\(v\), so \(w\forces B\to C\).

\paragraph{Cut.}
The premises are \(\Gamma_0\Rightarrow B\) and
\(B,\Pi_0\Rightarrow C\). At a world forcing
\(\Gamma_0,\Pi_0\), the first premise gives \(B\), and the second
then gives \(C\). This verifies precisely the displayed cut occurrence.

\paragraph{Mix.}
The premises are \(\Gamma_0\Rightarrow B\) and
\(B^n,\Pi_0\Rightarrow C\), where \(n\geq1\) is finite. From
forcing \(\Gamma_0,\Pi_0\), the first premise gives \(B\) at
\(w\). That same fact verifies each of the exactly \(n\) antecedent
occurrences in the second premise. It therefore gives \(C\) at
\(w\). There is no additional succedent formula or empty-succedent
case in this rule.

These are all the propositional rules. Exchange needs no further case
because a multiset has no presentation order. Finally, if an arbitrary
\(\Gamma\) derives \(A\), choose its finite witnessing context
\(\Gamma_0\). Every world forcing all of \(\Gamma\) forces each
occurrence of \(\Gamma_0\), and therefore forces \(A\). All cases
used only reflexivity, transitivity, and valuation persistence, so they
hold for preorders as well as partial orders.
\end{proof}
\end{annexitem}

\begin{annexitem}{annex:p8.ipl.complete.finite-support}{thm:p08-ipl-finite-support-operations}{Derivation operations with finite multiset supports}
\begin{proof}
Membership gives derivability by the identity sequent. Monotonicity follows
because a witnessing context over \(\Gamma\) is also a context over any
larger set. Finite character is part of the definition: the finitely many
formulas occurring in the witnessing antecedent form a finite subset of
\(\Gamma\).

\paragraph{Replacing derived premises.}
Suppose \(\Delta\derives_{\mathsf{IPL}}C\), witnessed by
\(D_1,\ldots,D_n\Rightarrow C\), where repeated formulas are allowed.
For each displayed occurrence choose a finite \(\Gamma\)-context
\(\Pi_i\) with a derivation of \(\Pi_i\Rightarrow D_i\).
There are only finitely many such choices. Cut this derivation against
the displayed occurrence of \(D_i\). Repeating for the original
\(n\) occurrences yields a derivation of
\(\Pi_1,\ldots,\Pi_n\Rightarrow C\); its antecedent is a finite
multiset of members of \(\Gamma\). If \(n=0\), the original
empty-antecedent derivation already witnesses the conclusion.
This proves substitution of derivations. Extensivity and monotonicity
of \(\operatorname{Cn}_{\mathsf{IPL}}\) follow from the preceding
paragraph. Substitution gives
\(\operatorname{Cn}_{\mathsf{IPL}}(\operatorname{Cn}_{\mathsf{IPL}}\Gamma)
\subseteq\operatorname{Cn}_{\mathsf{IPL}}\Gamma\); extensivity gives
the reverse inclusion, establishing idempotence.

\paragraph{Deduction.}
Suppose \(\Gamma\cup\{A\}\) derives \(B\). In a witnessing
context, group all occurrences equal to \(A\) and call the remainder
\(\Gamma_0\). Every member of \(\Gamma_0\) belongs to \(\Gamma\).
If there is at least one displayed \(A\), repeated CL leaves exactly
one. If there is none, WL adds one. Thus
\(A,\Gamma_0\Rightarrow B\) is derivable, and \(\to R\) gives
\(\Gamma_0\Rightarrow A\to B\). This argument also applies when
\(A\) itself belongs to \(\Gamma\).

Conversely, \(A\Rightarrow A\) and \(B\Rightarrow B\), followed
by \(\to L\), give \(A\to B,A\Rightarrow B\). Cut a finite
\(\Gamma\)-derivation of \(A\to B\) into this sequent. Its result
has a finite antecedent over \(\Gamma\cup\{A\}\), as required.

\paragraph{The connective operations.}
The same sequent \(A\to B,A\Rightarrow B\), followed by two
substitutions of derivations, proves modus ponens. For conjunction
introduction, let \(\Gamma_0\Rightarrow A\) and
\(\Delta_0\Rightarrow B\) be witnesses. WL extends each to the
same multiset antecedent \(\Gamma_0,\Delta_0\); \(\land R\)
then concludes \(A\land B\). For the first projection, identity
and WL give \(A,B\Rightarrow A\), and \(\land L\) gives
\(A\land B\Rightarrow A\). Starting from \(B\Rightarrow B\)
gives the other projection in the same explicit way:
\(A,B\Rightarrow B\), then \(A\land B\Rightarrow B\).
Cut either projection against a derivation of the conjunction.
The rules \(\lor R_1\) and \(\lor R_2\) give the two disjunction
introductions directly from a witnessing sequent. The initial sequent
\(\bottom\Rightarrow C\), followed by Cut, gives explosion.

For reasoning by cases, normalize the two branch witnesses by CL or WL,
exactly as in the deduction argument, to
\(A,\Gamma_0\Rightarrow C\) and
\(B,\Delta_0\Rightarrow C\). All formulas in the subscripted
contexts belong to \(\Gamma\). Use WL to give them the common
remaining antecedent \(\Gamma_0,\Delta_0\), then apply
\(\lor L\). The resulting sequent is
\(A\lor B,\Gamma_0,\Delta_0\Rightarrow C\).
Cut a finite \(\Gamma\)-derivation of \(A\lor B\) against the
displayed disjunction. This proves the assertion with a finite
\(\Gamma\)-support. All uses of Cut here are rules already present
on the sheet; no cut-elimination theorem is assumed.
\end{proof}
\end{annexitem}

\begin{annexitem}{annex:p8.ipl.complete.maximal-prime}{thm:p08-ipl-maximal-is-prime}{Why maximal avoidance is deductively closed, consistent, and prime}
\begin{proof}
First let \(S\) avoid \(C\) and derive \(A\lor B\). If both
\(S\cup\{A\}\) and \(S\cup\{B\}\) derived \(C\), reasoning
by cases would give \(S\derives_{\mathsf{IPL}}C\), a contradiction.
Classically, therefore, at least one extension avoids \(C\).

Now suppose \(T\) is maximal avoiding \(C\). If
\(T\derives_{\mathsf{IPL}}D\), adjoining \(D\) cannot create a
derivation of \(C\): substitution of the given derivation of \(D\)
would turn such a derivation into one from \(T\). Hence
\(T\cup\{D\}\) avoids \(C\), and maximality implies \(D\in T\).
This proves deductive closure. Membership of \(C\) itself would give
its derivability by identity, so \(C\notin T\). If \(T\) derived
\(\bottom\), explosion would give \(C\), proving consistency.

Finally, let \(A\lor B\in T\). It is derivable from \(T\) by
identity. The first paragraph gives a disjunct whose addition still
avoids \(C\); maximality places that disjunct in \(T\). Thus the
theory is prime. Conversely, for any deductively closed set,
\(A\in T\) implies \(T\derives A\) by identity and
\(T\derives A\) implies \(A\in T\) by closure, justifying the
membership/derivability equivalence used in the definition.
\end{proof}
\end{annexitem}

\begin{annexitem}{annex:p8.ipl.complete.prime-extension}{thm:p08-ipl-prime-extension}{Prime extension with arbitrary set cardinality}
\begin{proof}
The formula universe \(\mathcal F\) is a set: its finite construction
stages are sets and their union over the natural numbers is a set.
In particular, its power set is available. Let
\[
 \mathcal Q=\{S\subseteq\mathcal F:\Gamma\subseteq S
          \text{ and }S\not\derives_{\mathsf{IPL}}C\},
\]
ordered by inclusion. It is nonempty because \(\Gamma\in\mathcal Q\).
Every nonempty chain \(\mathcal C\subseteq\mathcal Q\) has upper
bound \(S^*=\bigcup\mathcal C\) in \(\mathcal Q\).
Indeed, it contains \(\Gamma\). If it derived \(C\), finitely
many antecedent formulas would witness this. Each belongs to some member
of \(\mathcal C\), and a finite collection of members of an inclusion
chain is contained in one of them. That member would derive \(C\),
contrary to its membership in \(\mathcal Q\). An empty witnessing
antecedent yields a derivation in every member of the nonempty chain
and gives the same contradiction. The empty chain has \(\Gamma\) as
an upper bound.

Zorn's lemma therefore gives a maximal member \(T\) of \(\mathcal Q\).
Every proper extension of \(T\) also contains \(\Gamma\), so it
cannot avoid \(C\) without contradicting this maximality. Thus \(T\)
is maximal avoiding \(C\). Lemma~\ref{thm:p08-ipl-maximal-is-prime}
makes it prime, with \(\Gamma\subseteq T\) and \(C\notin T\).
The proof used finite character, not a bound on the number of formulas
or premises. Substituting \(\bottom\) for \(C\) proves the final
assertion about arbitrary consistent premise sets.
\end{proof}
\end{annexitem}

\begin{annexitem}{annex:p8.ipl.complete.enumerated-extension}{thm:p08-ipl-enumerated-extension}{A choice-free construction for an enumerated language}
\begin{proof}
The recurrence defines a unique sequence in classical ZF: at each stage
the displayed derivability assertion or its negation determines the
next set. This uses classical excluded middle; it makes no assertion
that the test is computable. Inductively, \(S_n\subseteq S_{n+1}\)
and every \(S_n\) avoids \(C\). The initial case is the hypothesis;
at a successor stage the construction either keeps an avoiding set or
adds the formula only under the assertion that avoidance survives.

The union \(T\) avoids \(C\). Otherwise a finite witnessing context
would lie in some \(S_n\), contradicting that stage's avoidance.
For the empty context, use any stage, for example \(S_0\).
If \(D\notin T\), choose an index \(n\) with \(D_n=D\).
The construction could not have added \(D\) at this index, since
then \(D\in T\). Hence \(S_n\cup\{D\}\) derives \(C\),
and monotonicity gives \(T\cup\{D\}\derives C\).
Every proper extension of \(T\) contains some such \(D\) and thus
derives \(C\). This proves maximal avoidance, and the preceding
maximality lemma gives primeness.

No simultaneous choices of witnesses were made in this construction:
recursion specifies each set uniquely. For each fixed missing \(D\),
its least enumeration index could have been used. An enumeration of an
at most countable atom set yields an enumeration of finite formulas by
their finite construction codes; an empty atom set still has the formula
\(\bottom\) and countably many finite constructions. The statement
assumes the formula enumeration itself, so no inference from an
unspecified cardinal comparison is needed.
\end{proof}
\end{annexitem}

\begin{annexitem}{annex:p8.ipl.complete.failed-implication}{thm:p08-ipl-failed-implication}{A future prime theory separating an implication's operands}
\begin{proof}
If \(T\cup\{A\}\) derived \(B\), the deduction equivalence would
give \(T\derives_{\mathsf{IPL}}A\to B\). Deductive closure of
\(T\) would put \(A\to B\) in \(T\), contrary to the hypothesis.
Thus \(T\cup\{A\}\) avoids \(B\). Apply the appropriate
prime-extension lemma with this premise set and target \(B\).
It gives a prime \(U\supseteq T\cup\{A\}\) omitting \(B\),
which has all the claimed properties. The target of this extension is
\(B\); it need not be the formula omitted by some earlier root theory.
\end{proof}
\end{annexitem}

\begin{annexitem}{annex:p8.ipl.complete.truth}{thm:p08-ipl-canonical-truth}{Canonical model and the full truth induction}
\begin{proof}
The world collection \(W_{T_0}\) is a subset of the power set of
\(\mathcal F\), hence is a set. It is nonempty because it contains
\(T_0\). Inclusion is reflexive, transitive, and antisymmetric;
\(T_0\) is contained in every world, so it is a root. If
\(T\in V_{T_0}(p)\) and \(T\subseteq U\), then \(p\in T\subseteq U\),
so \(U\in V_{T_0}(p)\). This proves valuation persistence and makes
the canonical structure a model of the stipulated partial-order kind.

For every world, derivability and membership coincide by identity and
deductive closure. We prove the forcing/membership equivalence by
structural induction on \(D\), simultaneously at all worlds.

\paragraph{Atom and bottom.}
For an atom, forcing is defined by membership in its valuation, and the
valuation is defined by membership in the theory. No world forces
\(\bottom\); no prime theory contains \(\bottom\), since identity
would contradict consistency. Thus the equivalence holds in both cases.

\paragraph{Conjunction.}
If \(B\land C\in T\), its two derivable projections and closure
give \(B,C\in T\). Conversely, if both belong to \(T\), identity,
conjunction introduction, and closure give \(B\land C\in T\).
By induction these membership conditions are respectively forcing of
\(B\) and \(C\) at \(T\), exactly the forcing clause for their
conjunction.

\paragraph{Disjunction.}
Primeness gives \(B\in T\) or \(C\in T\) from \(B\lor C\in T\).
Conversely, either disjunct in \(T\) gives \(B\lor C\in T\) by
its introduction rule and closure. Induction and the local disjunction
clause prove the equivalence.

\paragraph{Implication: membership implies forcing.}
Suppose \(B\to C\in T\), and let \(U\in W_{T_0}\) extend \(T\)
with \(\mathcal M_{T_0},U\forces B\). By induction \(B\in U\).
Inclusion puts \(B\to C\) in \(U\); modus ponens and closure
therefore put \(C\) in \(U\). Induction now gives
\(\mathcal M_{T_0},U\forces C\). This holds for every accessible
\(U\), proving \(T\forces B\to C\).

\paragraph{Implication: failed membership implies failed forcing.}
Suppose \(B\to C\notin T\). The failed-implication extension lemma
gives a prime theory \(U\supseteq T\) with \(B\in U\) and
\(C\notin U\). Since \(T_0\subseteq T\subseteq U\), this
\(U\) is a world of the same canonical cone. The induction hypotheses
give \(U\forces B\) and \(U\not\forces C\). Thus \(U\) is an
accessible witness showing \(T\not\forces B\to C\).
Classical contraposition supplies the remaining direction of the
biconditional. The proof did not restrict \(U\) to omit any designated
root target; such a restriction would not justify this step.

This completes the truth induction. Finally, if \(T\forces D\) and
\(T\subseteq U\), the truth lemma gives \(D\in T\subseteq U\)
and then \(U\forces D\), establishing canonical persistence also
directly from membership.
\end{proof}
\end{annexitem}

\begin{annexitem}{annex:p8.ipl.complete.strong}{thm:p08-ipl-strong-completeness-full}{Strong completeness with an arbitrary premise set}
\begin{proof}
Soundness was proved for every propositional LJ rule and arbitrary
premise sets. For the converse, suppose
\(\Gamma\not\derives_{\mathsf{IPL}}A\). Prime extension gives
\(T_0\supseteq\Gamma\) with \(A\notin T_0\). Form the canonical
model above \(T_0\). Its truth lemma says that the root forces every
member of \(\Gamma\), because each belongs to \(T_0\), and that
the root does not force \(A\). This is a rooted partially ordered
countermodel to the proposed consequence, so classical contraposition
proves completeness.

The same single countermodel belongs to the classes of all partial
orders and all preorders. Soundness holds in all those classes, so the
three separately specified consequences have precisely the same
derivable sequents. The construction used no conjunction of the
premises and no enumeration of \(\Gamma\). An infinite premise set
is handled by containment in \(T_0\) and the pointwise truth lemma.
For a supplied formula enumeration, replace each use of the Zorn
extension lemma by the enumerated one. All remaining arguments are
classical ZF arguments, which proves the final metatheoretic assertion.
\end{proof}
\end{annexitem}

\begin{annexitem}{annex:p8.ipl.complete.compactness}{thm:p08-ipl-semantic-compactness}{Finite support and satisfiability compactness}
\begin{proof}
If \(\Gamma\) entails \(A\), completeness gives a derivation with
finite antecedent \(\Gamma_0\). Let \(\Gamma_f\) be the finite
set of formulas occurring there. The same derivation witnesses
\(\Gamma_f\derives A\), and soundness gives \(\Gamma_f\entails A\).

If a world forces \(\Gamma\), it forces each finite subset.
Conversely, suppose every finite subset is forced somewhere. Were
\(\Gamma\derives\bottom\), a finite supporting subset would
derive \(\bottom\). Soundness at its assumed satisfying world would
force \(\bottom\), which no world can force. Hence
\(\Gamma\not\derives\bottom\). Apply strong completeness's
countermodel construction with target \(\bottom\). The resulting
rooted partial-order model has a root forcing all of \(\Gamma\).
The witnesses for different finite subsets need not be related, and no
choices of a simultaneous family of these witnesses are used.
\end{proof}
\end{annexitem}

\begin{annexitem}{annex:p8.ipl.complete.filtration}{thm:p08-ipl-filtration-truth}{Finite filtration with the implication clauses verified}
\begin{proof}
Equality of the finite truth sets is an equivalence relation. Changing
representatives in either argument of \(\leq_\Sigma\) leaves those
sets unchanged, so the relation is well-defined. It is reflexive and
transitive because inclusion is. If two classes are related in both
directions, their truth sets are equal and the classes themselves are
equal. Thus the quotient order is antisymmetric. Mapping a class to its
truth set injects \(W_\Sigma\) into \(\mathcal P(\Sigma)\),
giving \(|W_\Sigma|\leq2^{|\Sigma|}\). The quotient is nonempty
because \(W\) is nonempty.

For \(p\in\Sigma\), equivalent representatives agree on forcing
\(p\), so they agree on membership in \(V(p)\). This makes
\(V_\Sigma(p)\) well-defined. If \([w]_\Sigma\leq_\Sigma[v]_\Sigma\)
and \([w]_\Sigma\in V_\Sigma(p)\), then
\(p\in\tau_\Sigma(w)\subseteq\tau_\Sigma(v)\), whence
\([v]_\Sigma\in V_\Sigma(p)\). The valuation for an atom outside
\(\Sigma\) is empty and is persistent vacuously. Thus the filtered
structure is a model for the entire language, not just a partial
assignment on its relevant atoms.

Original accessibility implies filtered accessibility: if \(w\leq v\),
formula persistence gives \(\tau_\Sigma(w)\subseteq\tau_\Sigma(v)\).
In particular, an original root \(r\) gives
\([r]_\Sigma\leq_\Sigma[w]_\Sigma\) for every class, proving
root preservation.

We now prove preservation of forcing for each \(D\in\Sigma\) by
structural induction. Subformula closure puts the operands of any
compound \(D\in\Sigma\) in \(\Sigma\), so the induction
hypotheses apply to them at every original and filtered world.

\paragraph{Atoms, bottom, conjunction, and disjunction.}
For an atom \(p\in\Sigma\), the valuation definition gives exactly
the desired equivalence. Neither model forces \(\bottom\).
For \(B\land C\), forcing in either model means forcing both
operands at that model's world; applying both induction equivalences
gives the assertion. For \(B\lor C\), forcing means forcing at
least one operand at that same world. The two induction equivalences
again give both directions.

\paragraph{Implication preserved from the original to the quotient.}
Suppose \(B\to C\in\Sigma\) and \(\mathcal M,w\forces B\to C\).
Take any \([v]_\Sigma\geq_\Sigma[w]_\Sigma\) with
\(\mathcal M_\Sigma,[v]_\Sigma\forces B\). Because the implication
formula itself is in \(\Sigma\), it belongs to
\(\tau_\Sigma(w)\) and hence to \(\tau_\Sigma(v)\). Therefore
\(\mathcal M,v\forces B\to C\), even if \(w\not\leq v\)
in the original frame. The induction hypothesis gives
\(\mathcal M,v\forces B\). Reflexivity at \(v\) and its forced
implication yield \(\mathcal M,v\forces C\), and induction gives
\(\mathcal M_\Sigma,[v]_\Sigma\forces C\). The filtered
implication clause now gives
\(\mathcal M_\Sigma,[w]_\Sigma\forces B\to C\).

\paragraph{Implication reflected from the quotient to the original.}
Suppose \(\mathcal M_\Sigma,[w]_\Sigma\forces B\to C\).
For any original \(v\geq w\) with \(\mathcal M,v\forces B\),
original accessibility gives filtered accessibility, and induction gives
\(\mathcal M_\Sigma,[v]_\Sigma\forces B\). The forced implication
at \([w]_\Sigma\) yields the filtered forcing of \(C\) at
\([v]_\Sigma\); induction reflects it to
\(\mathcal M,v\forces C\). This proves the original implication
clause at \(w\). Equivalently, an original future world forcing
\(B\) and failing \(C\) retains both properties after filtration
and witnesses failure of the filtered implication. Both directions are
therefore proved without assuming that every filtered comparison lifts
to an original comparison.
\end{proof}
\end{annexitem}

\begin{annexitem}{annex:p8.ipl.complete.finite-model}{thm:p08-ipl-finite-rooted-countermodel}{The finite rooted countermodel promised by IPL adequacy}
\begin{proof}
Strong completeness supplies a canonical partial-order model with root
\(T_0\) forcing every member of \(\Gamma\) and not forcing \(A\).
Each formula is a finite construction, and \(\Gamma\cup\{A\}\)
is finite. Its full subformula set \(\Sigma\) is therefore finite
and is closed under subformulas. Filter the model through \(\Sigma\).
The filtration lemma makes the quotient a partial order with at most
\(2^{|\Sigma|}\) worlds, preserves its root, and preserves forcing
of every formula in \(\Gamma\cup\{A\}\). Thus the quotient root
forces all the premises and fails the conclusion, as claimed. This
argument also covers \(\Gamma=\varnothing\). The cardinal bound
depends on the finite input formulas, not on the cardinality of the
ambient atom set or of the canonical frame.
\end{proof}
\end{annexitem}

\begin{annexitem}{annex:p8.ipl.complete.decision}{thm:p08-ipl-finite-input-decision}{Termination and correctness of finite model search}
\begin{proof}
Expand negation, compute the finite subformula set \(\Sigma\), and
let \(m=|\Sigma|\). Let \(P_0\) be the finite set of atoms in
the input. For each integer \(n\) from \(1\) through \(2^m\),
use world set \(\{0,\ldots,n-1\}\). Enumerate all its binary
relations and retain precisely the reflexive, transitive, antisymmetric
ones. For each retained relation enumerate each possible root and keep
it if it lies below every world. For each such rooted frame enumerate
all assignments of upward-closed subsets to the atoms in \(P_0\).
Assign the empty set to every other atom.

Each list is finite: there are \(2^{n^2}\) binary relations, \(n\)
possible roots, and at most \(2^{n|P_0|}\) atom assignments. The
frame properties and upward-closure conditions are finite checks.
Forcing of each input formula is computed by structural recursion;
the implication clause checks finitely many future worlds. If any root
forces the premises and fails the conclusion, output nonderivable.
Soundness proves the correctness of that output.

If the search finishes with no countermodel, output derivable. Were the
conclusion not derivable, the finite-countermodel theorem would provide
a rooted partial-order countermodel with at most \(2^m\) worlds.
Relabel its finite world set as \(\{0,\ldots,n-1\}\).
The search includes that frame and the restrictions of its valuations
to \(P_0\). Setting the other atoms false cannot change any input
formula, as structural induction on these formulas shows: their atomic
cases all lie in \(P_0\), and every connective clause uses only its
operands and the unchanged frame. The search would consequently have
found the countermodel, a contradiction. The positive output is
therefore also correct, and the finite search always terminates.

An explicit input encoding with decidable atom equality suffices to
compute \(P_0\), \(\Sigma\), and all the stated finite lists.
No algorithm for testing membership in an arbitrary infinite premise
set is assumed or furnished by this finite-input result.
\end{proof}
\end{annexitem}
